\chapter{Options et arithmétique obligataire à toute vitesse}\label{ch:iv2:options-and-bond-arithmetic-at-speed}

\begin{ivtrains}
Les ordres de grandeur qu'un recruteur s'attend à ce que l'on connaisse de tête~: valeurs d'options à la monnaie et grecques sans calculatrice, parité à la main, DV01 et variations de rendement, carry, fair value des ETF et des contrats à terme. On donne d'abord l'estimation, puis l'écart qu'elle peut avoir. Les taux sont simples, sauf indication contraire dans l'énoncé.
\end{ivtrains}

\begin{question}[$\star$ \iqroles{QT, Trader, QR} \iqfirm{teneur de marché, hedge fund, banque} \iqlevel{junior}]\label{iq:iv2:options-and-bond-arithmetic-at-speed:1}
Une action se traite à 100~\$ avec une volatilité de 20~\% par an, et les taux d'intérêt sont nuls. Sans calculatrice, combien vaut un call à la monnaie d'un an~? Et un call à la monnaie de six mois sur une action à 50~\$ avec une volatilité de 30~\%~? Donnez la règle que vous utilisez et dites à peu près à quel point elle est précise.
\end{question}

\begin{question}[$\star$ \iqroles{QT, Trader} \iqfirm{teneur de marché, hedge fund} \iqlevel{junior}]\label{iq:iv2:options-and-bond-arithmetic-at-speed:2}
Une action se traite à 200~\$ avec une volatilité implicite de 40~\%. Les taux sont nuls. Combien coûte un straddle à la monnaie de trois mois, soit un call et un put au même strike~? De combien l'action doit-elle bouger pour que le straddle soit à l'équilibre à l'échéance~? Quelle est la variation quotidienne typique, en pourcentage, avec 252 jours de bourse dans l'année~?
\end{question}

\begin{question}[$\star$ \iqroles{QT, Trader, QR} \iqfirm{teneur de marché, hedge fund} \iqlevel{junior}]\label{iq:iv2:options-and-bond-arithmetic-at-speed:3}
Les taux sont nuls. Quel est le delta d'un call à la monnaie d'un an sur une action de volatilité 20~\%, dans le modèle de Black--Scholes~? Et celui du put de même strike~? Estimez-les de tête, et dites pourquoi le delta du call n'est pas exactement égal à un demi.
\end{question}

\begin{question}[$\star$ \iqroles{QT, Trader} \iqfirm{teneur de marché, hedge fund, banque} \iqlevel{junior}]\label{iq:iv2:options-and-bond-arithmetic-at-speed:4}
Une action à 100~\$ dispose d'options à la monnaie. De combien varie la valeur d'un call à la monnaie d'un an lorsque la volatilité implicite augmente d'un point de pourcentage~? Et pour un call de trois mois~? Et pour un call d'un an sur une action à 500~\$~? Estimez chacun de tête, et dites de quoi dépend ce chiffre.
\end{question}

\begin{question}[$\star$ \iqroles{QT, Trader, QR} \iqfirm{teneur de marché, hedge fund, banque} \iqlevel{junior}]\label{iq:iv2:options-and-bond-arithmetic-at-speed:5}
Une action se traite à 100~\$ et ne verse aucun dividende. Le taux d'intérêt d'un an est de 5~\%, capitalisé en continu. Un call d'un an de strike 100~\$ coûte 12~\$. Combien doit coûter le put d'un an de même strike~? Donnez une estimation rapide, puis la valeur exacte.
\end{question}

\begin{question}[$\star$ \iqroles{QT, Trader, Risque} \iqfirm{banque, hedge fund} \iqlevel{junior}]\label{iq:iv2:options-and-bond-arithmetic-at-speed:6}
Vous détenez 25~millions de dollars d'une obligation de duration modifiée 6.5 et de convexité 60. Les rendements montent de 30 points de base, soit 0.30 point de pourcentage. Quelle est votre perte due à la duration seule~? Que devient cette perte en tenant compte de la convexité~?
\end{question}

\begin{question}[$\star$ \iqroles{QT, Trader, Risque} \iqfirm{banque, hedge fund} \iqlevel{junior}]\label{iq:iv2:options-and-bond-arithmetic-at-speed:7}
Vous détenez une obligation d'une valeur de marché de 10~millions de dollars, de duration modifiée 8.2. Quel est son DV01, soit la variation de valeur pour un mouvement d'un point de base du rendement~? Un contrat à terme sur obligations a un DV01 de 65~\$ par contrat. Combien de contrats vendez-vous pour la couvrir~?
\end{question}

\begin{question}[$\star$ \iqroles{QT, Trader, QR} \iqfirm{banque, hedge fund} \iqlevel{junior}]\label{iq:iv2:options-and-bond-arithmetic-at-speed:8}
Une obligation zéro-coupon rembourse 100 dans 10 ans, et son rendement est de 4~\%, capitalisé annuellement. Quel est son prix, de tête~? Quelles sont ses durations de Macaulay et modifiée~? Quelle est, à peu près, la variation de son prix si son rendement monte à 4.1~\%~?
\end{question}

\begin{question}[$\star$ \iqroles{QT, Trader, QR} \iqfirm{banque, hedge fund} \iqlevel{junior}]\label{iq:iv2:options-and-bond-arithmetic-at-speed:9}
Le taux zéro d'un an est de 3~\%, et le taux zéro de deux ans de 4~\%, les deux capitalisés annuellement. Quel est le taux forward d'un an, démarrant dans un an~? Donnez la réponse rapide et la réponse exacte, et dites pourquoi la réponse rapide fonctionne.
\end{question}

\begin{question}[$\star$ \iqroles{QT, Trader} \iqfirm{banque, hedge fund, teneur de marché} \iqlevel{junior}]\label{iq:iv2:options-and-bond-arithmetic-at-speed:10}
Le spot est de 1.10 dollar pour un euro. Le taux en dollars à un an est de 5~\%, et le taux en euros à un an de 3~\%, les deux exprimés en taux simples. Quel est le cours forward à un an, en approximation rapide puis exacte~? Si quelqu'un cote un forward de 1.13, comment verrouillez-vous un profit~?
\end{question}

\begin{question}[$\star$ \iqroles{QT, Trader} \iqfirm{teneur de marché, hedge fund} \iqlevel{junior}]\label{iq:iv2:options-and-bond-arithmetic-at-speed:11}
Les parts d'un ETF sont créées par unités de 50\,000. Le panier qui sous-tend une part d'ETF coûte 49.90~\$ à l'achat. Acheter le panier et créer coûte en tout 0.03~\$ par part d'ETF. L'ETF est coté avec un cours acheteur de 50.02~\$ et un cours vendeur de 50.04~\$. Quelle transaction faites-vous, et que gagnez-vous par unité de création~?
\end{question}

\begin{question}[$\star$ \iqroles{QT, Trader} \iqfirm{teneur de marché, hedge fund} \iqlevel{junior}]\label{iq:iv2:options-and-bond-arithmetic-at-speed:12}
Le straddle à la monnaie de trois mois d'une action à 100~\$ se traite à 12~\$, et les taux sont nuls. Quelle volatilité implicite cela suppose-t-il, de tête~? Quel est le mouvement d'un écart-type à l'échéance, en dollars~?
\end{question}

\begin{question}[$\star\star$ \iqroles{QT, Trader, QR} \iqfirm{teneur de marché, hedge fund} \iqlevel{junior}]\label{iq:iv2:options-and-bond-arithmetic-at-speed:13}
Les taux sont nuls. Un call à la monnaie d'un an sur une action à 100~\$ de volatilité 20~\% vaut environ 8~\$. De tête, estimez son gamma, soit la variation du delta pour une variation de 1~\$ du cours de l'action. Estimez son theta d'un jour, avec 252 jours de bourse dans l'année, de deux façons. De combien l'action doit-elle bouger en un jour pour que le gain du gamma compense ce theta~?
\end{question}

\begin{question}[$\star\star$ \iqroles{QT, Trader} \iqfirm{teneur de marché, hedge fund, banque} \iqlevel{junior}]\label{iq:iv2:options-and-bond-arithmetic-at-speed:14}
Une action se traite à 100~\$ et verse un dividende de 2~\$ dans trois mois. Le taux d'intérêt est de 4~\% par an, capitalisé en continu. Les options européennes de six mois ont un strike de 95~\$. Quelle est la différence entre le prix du call et celui du put, de tête~? Quelle serait-elle si vous aviez oublié le dividende~?
\end{question}

\begin{question}[$\star\star$ \iqroles{QT, Trader} \iqfirm{teneur de marché, hedge fund} \iqlevel{junior}]\label{iq:iv2:options-and-bond-arithmetic-at-speed:15}
Une action se traite à 100~\$, et les options européennes de six mois de strike 100~\$ ont un prix de call de 5.50~\$ et un prix de put de 5.00~\$. Le taux d'intérêt de six mois est de 3~\% par an, en taux simple. Quel prix forward les options impliquent-elles~? Que donne le seul carry~? Quel est l'écart, et que pourrait-il représenter~?
\end{question}

\begin{question}[$\star\star$ \iqroles{QT, Trader, Risque} \iqfirm{banque, hedge fund} \iqlevel{junior}]\label{iq:iv2:options-and-bond-arithmetic-at-speed:16}
Une obligation de 10 ans verse un coupon annuel de 5~\% et son rendement est de 5~\%. Pourquoi son prix vaut-il 100, et quelle formule rapide donne sa duration modifiée~? Estimez son prix si son rendement monte à 5.1~\%, puis comparez-le au prix exact.
\end{question}

\begin{question}[$\star\star$ \iqroles{QT, Trader} \iqfirm{banque, hedge fund} \iqlevel{junior, confirmé}]\label{iq:iv2:options-and-bond-arithmetic-at-speed:17}
Une obligation de 5 ans a un rendement de 4.5~\% avec une duration modifiée de 4.5. Vous la financez à 3.5~\% sur le marché du repo. La courbe des taux autour d'elle est croissante, de 10 points de base par année de maturité, et elle reste inchangée pendant les trois mois suivants. Quels sont votre carry, votre roll-down et votre total, en points de prix, sur les trois mois~? De combien son rendement peut-il monter avant que vous ne perdiez de l'argent~?
\end{question}

\begin{question}[$\star\star$ \iqroles{QT, Trader} \iqfirm{teneur de marché, hedge fund} \iqlevel{junior, confirmé}]\label{iq:iv2:options-and-bond-arithmetic-at-speed:18}
Un ETF ne détient que des actions japonaises et se traite aux États-Unis pendant que Tokyo est fermée. Sa valeur liquidative publiée, issue du dernier cours de clôture à Tokyo, est de 40.00~\$. Depuis cette clôture, les contrats à terme sur le même indice ont progressé de 1~\%, et la devise n'a pas bougé. L'ETF se traite à 40.50~\$. Quelle prime la valeur publiée affiche-t-elle~? Quelle est la prime réelle~? Votre coût aller-retour est de 20 points de base~: y a-t-il une transaction à réaliser~?
\end{question}

\begin{question}[$\star\star$ \iqroles{QT, Trader} \iqfirm{teneur de marché, hedge fund, banque} \iqlevel{junior}]\label{iq:iv2:options-and-bond-arithmetic-at-speed:19}
Un indice boursier vaut 4\,000. Son contrat à terme de trois mois expire exactement dans un quart d'année. Les taux d'intérêt sont de 5~\% par an, le rendement du dividende est de 2~\% par an, et vous utilisez le carry simple. Quelle est la fair value~? Le contrat à terme se traite à 4\,036. Quelle transaction faites-vous, combien gagnez-vous en points de base du spot, et que reste-t-il si l'entrée et la sortie vous coûtent ensemble 10 points de base~?
\end{question}

\begin{question}[$\star\star$ \iqroles{QT, Trader, QR} \iqfirm{teneur de marché, hedge fund} \iqlevel{junior, confirmé}]\label{iq:iv2:options-and-bond-arithmetic-at-speed:20}
Les taux sont nuls. Une action se traite à 100~\$ avec une volatilité de 20~\%. Un call spread d'un an, achat du call 100 et vente du call 110, coûte environ 3.67~\$. Que vous apprend ce prix sur la probabilité de terminer au-dessus d'environ 105~\$, et pourquoi~? Que donne le modèle de Black--Scholes pour cette probabilité~?
\end{question}

\begin{question}[$\star\star$ \iqroles{QT, Trader, QR} \iqfirm{teneur de marché, hedge fund} \iqlevel{junior, confirmé}]\label{iq:iv2:options-and-bond-arithmetic-at-speed:21}
Les taux sont nuls. Une action se traite à 100~\$. Un call d'un an de strike 110~\$ a une volatilité implicite de 20~\%. Traitez le mouvement de l'action comme une loi normale, d'écart-type 20~\$, et utilisez $\Phi(0.5) \approx 0.69$ et une densité normale d'environ 0.35 en 0.5. Quelle valeur estimez-vous~? Comment se compare-t-elle à la valeur de Black--Scholes, et pourquoi est-elle plus basse~?
\end{question}

\begin{question}[$\star\star$ \iqroles{QT, QR, Trader} \iqfirm{teneur de marché, hedge fund} \iqlevel{junior, confirmé}]\label{iq:iv2:options-and-bond-arithmetic-at-speed:22}
Les taux sont nuls. Une action se traite à 100~\$. Quel delta donne la règle empirique pour un call dont le strike est situé à un écart-type au-dessus de l'action, l'écart-type étant $\sigma\sqrt{T}$ exprimé en pourcentage~? Quel est le delta de Black--Scholes si l'option a un an et une volatilité de 20~\%, avec un strike de 120~? Et si elle a une durée de vie courte, avec une volatilité totale de 5~\% et un strike de 105~? Pourquoi ces deux valeurs diffèrent-elles~?
\end{question}

\begin{question}[$\star\star$ \iqroles{QT, Trader, Risque} \iqfirm{banque, hedge fund} \iqlevel{junior, confirmé}]\label{iq:iv2:options-and-bond-arithmetic-at-speed:23}
Pour 1~million de dollars de valeur nominale, les DV01 des obligations de 2 ans, 5 ans et 10 ans valent 190~\$, 460~\$ et 830~\$. Vous achetez 10~millions de dollars de l'obligation à 5 ans et vendez les deux autres, avec la moitié du risque sur chaque aile et un DV01 net nul. Combien vendez-vous de chacune~? Les rendements sont de 4.0~\%, 4.4~\% et 4.6~\%. Quel est le spread du butterfly, soit deux fois le rendement du ventre moins celui des ailes~? Que gagnez-vous ou perdez-vous si seul le rendement à 5 ans baisse de 5 points de base, et si les trois montent de 10~?
\end{question}

\begin{question}[$\star\star\star$ \iqroles{QT, Trader, QR} \iqfirm{teneur de marché, hedge fund} \iqlevel{junior, confirmé}]\label{iq:iv2:options-and-bond-arithmetic-at-speed:24}
Les taux sont nuls. Vous vendez un straddle à la monnaie d'un mois sur une action à 100~\$, à une volatilité implicite de 30~\%, et vous la couvrez souvent en delta, en utilisant les deltas calculés à 30~\%. L'action fluctue ensuite d'environ 1~\% par jour, sans tendance, avec 252 jours de bourse dans l'année. Qu'encaissez-vous~? Quelle volatilité représente un mouvement de 1~\% par jour~? Quel est votre profit espéré, et pourquoi ce n'est pas simplement le vega multiplié par l'écart de volatilité, à gamma constant~?
\end{question}

\begin{question}[$\star\star\star$ \iqroles{QT, Trader} \iqfirm{teneur de marché, hedge fund} \iqlevel{junior, confirmé}]\label{iq:iv2:options-and-bond-arithmetic-at-speed:25}
Les taux sont nuls. Une action à 100~\$ publie ses résultats dans cinq jours de bourse. Le straddle à la monnaie qui expire juste après la publication coûte 8~\$, et les jours ordinaires l'action a une volatilité de 30~\%, avec 252 jours de bourse dans l'année. Quel écart-type total le straddle suppose-t-il~? Quel mouvement cela suppose-t-il pour le seul jour de publication~? Que feriez-vous si vous pensiez que le mouvement lié aux résultats serait d'environ 6~\%~?
\end{question}
